% ============================================================
% 培优专题9 — 数列递推与不等式
% 编译命令：xelatex -interaction=nonstopmode 培优专题9_数列递推与不等式.tex
% ============================================================
\documentclass[12pt, a4paper]{ctexart}
\usepackage[margin=2.2cm]{geometry}
\usepackage{amsmath, amssymb}
\usepackage{enumitem}
\usepackage{tcolorbox}
\usepackage{booktabs}
\usepackage{array}
\usepackage{fancyhdr}
\usepackage{xcolor}
\usepackage{tikz}
\usepackage{hyperref}
\usepackage{graphicx}
\usepackage{multirow}
\usepackage{pgfplots}
\pgfplotsset{compat=1.18}

% 颜色定义
\definecolor{mainblue}{RGB}{41, 98, 168}
\definecolor{lightblue}{RGB}{230, 240, 255}
\definecolor{tipgreen}{RGB}{34, 139, 87}
\definecolor{lightgreen}{RGB}{235, 250, 240}
\definecolor{warnred}{RGB}{200, 50, 50}
\definecolor{lightred}{RGB}{255, 240, 240}
\definecolor{examorange}{RGB}{210, 120, 20}
\definecolor{lightorange}{RGB}{255, 248, 235}
\definecolor{apppurple}{RGB}{120, 60, 170}
\definecolor{lightpurple}{RGB}{245, 238, 255}

% 页眉页脚
\pagestyle{fancy}
\fancyhf{}
\fancyhead[L]{\small\textcolor{mainblue}{\textbf{高中数学培优 · 数列递推与不等式}}}
\fancyhead[R]{\small\textcolor{mainblue}{培优专题9}}
\fancyfoot[C]{\small —\ \thepage\ —}
\renewcommand{\headrulewidth}{0.6pt}
\renewcommand{\headrule}{\hbox to\headwidth{\color{mainblue}\leaders\hrule height \headrulewidth\hfill}}
\renewcommand{\footrulewidth}{0pt}

% tcolorbox 样式
\newtcolorbox{knowledgebox}[1][]{%
  colback=lightblue, colframe=mainblue,
  fonttitle=\bfseries\large, title=#1,
  arc=3mm, boxrule=0.8pt,
  left=8pt, right=8pt, top=6pt, bottom=6pt
}

\newtcolorbox{tipbox}[1][]{%
  colback=lightgreen, colframe=tipgreen,
  fonttitle=\bfseries, title=#1,
  arc=2mm, boxrule=0.6pt,
  left=6pt, right=6pt, top=4pt, bottom=4pt
}

\newtcolorbox{warnbox}[1][]{%
  colback=lightred, colframe=warnred,
  fonttitle=\bfseries, title=#1,
  arc=2mm, boxrule=0.6pt,
  left=6pt, right=6pt, top=4pt, bottom=4pt
}

\newtcolorbox{examplebox}[1][]{%
  colback=lightorange, colframe=examorange,
  fonttitle=\bfseries, title=#1,
  arc=3mm, boxrule=0.8pt,
  left=8pt, right=8pt, top=6pt, bottom=6pt
}

\newtcolorbox{appbox}[1][]{%
  colback=lightpurple, colframe=apppurple,
  fonttitle=\bfseries, title=#1,
  arc=3mm, boxrule=0.8pt,
  left=8pt, right=8pt, top=6pt, bottom=6pt
}

\setCJKmainfont{Hiragino Sans GB}
\setCJKsansfont{Heiti SC}

\begin{document}

% ============ 标题页 ============
\begin{center}
{\Huge\bfseries\textcolor{mainblue}{培优专题9\quad 数列递推与不等式}}\\[12pt]
{\Large\textcolor{gray}{培\quad 优\quad 讲\quad 义}}\\[8pt]
\rule{0.6\textwidth}{1pt}\\[8pt]
{\normalsize 适用：高二高三拔高\quad |\quad 课时：3课时}
\end{center}
\vspace{10pt}

% ============ 第一部分：知识点讲解 ============
\section*{\textcolor{mainblue}{一、知识点讲解}}

% ------ 1.1 特征方程法 ------
\subsection*{1.1\quad 特征方程法}

\begin{knowledgebox}[特征方程法基本原理]
对于二阶线性递推数列
\[
a_{n+2} = p\,a_{n+1} + q\,a_n \quad (p,q\text{为常数},\;q\neq 0)
\]
令 $r^2 = pr + q$，称此方程为\textbf{特征方程}。

设其两根为 $r_1, r_2$：

\textbf{情形一：}$r_1 \neq r_2$（两个不等实根）
\[
a_n = A\,r_1^{\,n} + B\,r_2^{\,n}
\]
其中 $A, B$ 由初始条件 $a_1, a_2$ 确定。

\textbf{情形二：}$r_1 = r_2 = r$（重根）
\[
a_n = (A + Bn)\,r^{\,n}
\]

\textbf{情形三：}$r_1, r_2$ 为共轭复根 $r = \alpha \pm \beta i$
\[
a_n = R^n\left(A\cos n\theta + B\sin n\theta\right)
\]
其中 $R = |r|,\;\theta = \arg(r)$。
\end{knowledgebox}

\begin{tipbox}[推广：一阶线性递推]
对于 $a_{n+1} = pa_n + q$（$p\neq 1$），其特征方程为 $r = pr + q$，解得不动点 $r = \dfrac{q}{1-p}$。令 $b_n = a_n - r$，则 $b_{n+1} = pb_n$，为等比数列。
\end{tipbox}

% ------ 1.2 不动点法 ------
\subsection*{1.2\quad 不动点法}

\begin{knowledgebox}[不动点法]
对于递推关系 $a_{n+1} = f(a_n)$，若方程 $f(x) = x$ 有实数解 $x_0$，则称 $x_0$ 为 $f$ 的\textbf{不动点}。

\textbf{类型一：}$a_{n+1} = pa_n + q$（$p\neq 1$）

不动点 $x_0 = \dfrac{q}{1-p}$，令 $b_n = a_n - x_0$，则 $b_{n+1} = pb_n$。

\textbf{类型二：}$a_{n+1} = \dfrac{a\,a_n + b}{c\,a_n + d}$（分式线性递推）

不动点为方程 $x = \dfrac{ax+b}{cx+d}$ 的根，即 $cx^2 + (d-a)x - b = 0$ 的根。

若有两个不等不动点 $\alpha, \beta$，令
\[
\frac{a_n - \alpha}{a_n - \beta} = t \cdot \frac{a_{n-1} - \alpha}{a_{n-1} - \beta}
\]
可化为等比数列求解。
\end{knowledgebox}

\begin{warnbox}[易错提醒]
使用不动点法时，需验证递推函数在不动点附近有意义。若不动点处 $f'(x_0) = 0$，则收敛速度会加快（二次收敛）。特别注意分式递推中分母不为零的限制。
\end{warnbox}

% ------ 1.3 递推数列周期性 ------
\subsection*{1.3\quad 递推数列周期性}

\begin{knowledgebox}[周期数列的判定]
数列 $\{a_n\}$ 若存在正整数 $T$，使得 $a_{n+T} = a_n$ 对一切 $n$ 成立，则称其为\textbf{周期数列}，$T$ 为一个周期。

\textbf{常见周期性递推模型：}

(1) 二阶循环：$a_{n+2} = \dfrac{1}{a_n}$，周期 $T=4$（一般情况）。

(2) 三角型：$a_{n+1} = 2a_n\sqrt{1-a_n^2}$，令 $a_n = \sin\theta_n$ 可化为 $\theta_{n+1} = 2\theta_n$。

(3) 取模型：$a_{n+1} \equiv f(a_n) \pmod{m}$，至多 $m$ 项后出现周期。

\textbf{判定技巧：}
\begin{itemize}[nosep]
\item 计算若干项，观察是否重复；
\item 利用函数迭代图分析；
\item 利用模运算的有限性。
\end{itemize}
\end{knowledgebox}

% ------ 1.4 数列不等式证明 ------
\subsection*{1.4\quad 数列不等式证明}

\begin{knowledgebox}[常用证明策略]
\textbf{策略一：数学归纳法}
\begin{itemize}[nosep]
\item 验证 $n=1$（或起始值）成立；
\item 假设 $n=k$ 成立，证明 $n=k+1$ 也成立。
\end{itemize}

\textbf{策略二：单调性论证}
\begin{itemize}[nosep]
\item 证明 $a_{n+1} - a_n > 0$（或 $< 0$）；
\item 利用 $a_n$ 与目标值的差构造单调性。
\end{itemize}

\textbf{策略三：放缩法}
\begin{itemize}[nosep]
\item 利用 $a_n < b_n$（或 $>$）逐层放缩；
\item 常配合裂项、均值不等式等工具。
\end{itemize}

\textbf{策略四：构造辅助数列}
\begin{itemize}[nosep]
\item 令 $b_n = a_n - c$ 或 $b_n = \dfrac{a_n}{c}$ 等，化简递推关系后证明。
\end{itemize}
\end{knowledgebox}

% ------ 1.5 裂项放缩 ------
\subsection*{1.5\quad 裂项放缩}

\begin{knowledgebox}[裂项放缩核心公式]
\textbf{基本裂项：}
\[
\frac{1}{n(n+1)} = \frac{1}{n} - \frac{1}{n+1}
\]
\[
\frac{1}{n(n+k)} = \frac{1}{k}\left(\frac{1}{n} - \frac{1}{n+k}\right)
\]

\textbf{常用放缩：}
\[
\frac{1}{n^2} < \frac{1}{n(n-1)} = \frac{1}{n-1} - \frac{1}{n} \quad (n\geq 2)
\]
\[
\frac{1}{n^2} > \frac{1}{n(n+1)} = \frac{1}{n} - \frac{1}{n+1}
\]
\[
\frac{1}{\sqrt{n} + \sqrt{n+1}} = \sqrt{n+1} - \sqrt{n}
\]

\textbf{进阶放缩：}
\[
\frac{1}{n^2} < \frac{1}{n^2 - \frac{1}{4}} = \frac{1}{\left(n-\frac{1}{2}\right)\left(n+\frac{1}{2}\right)} = \frac{1}{2n-1} - \frac{1}{2n+1} \cdot \frac{1}{2}
\]
更精确地：
\[
\frac{1}{n^2} < \frac{1}{n^2 - \frac{1}{4}} = \frac{1}{\left(n-\frac{1}{2}\right)\left(n+\frac{1}{2}\right)}
\]
\end{knowledgebox}

\begin{tipbox}[放缩原则]
放缩方向要一致（同向传递），且放缩程度要"恰到好处"——既不过大（导致目标无法达到），也不过小（导致无法求和）。一般先试简单放缩，不成功再调整。
\end{tipbox}

% ============ 第二部分：例题精讲 ============
\section*{\textcolor{examorange}{二、例题精讲}}

% ------ 例题1 ------
\begin{examplebox}[【例题1】特征方程法求通项]
已知数列 $\{a_n\}$ 满足 $a_1 = 1,\;a_2 = 3$，且 $a_{n+2} = 3a_{n+1} - 2a_n$，求通项公式。
\end{examplebox}

\textbf{解：}特征方程为 $r^2 = 3r - 2$，即 $r^2 - 3r + 2 = 0$。

解得 $r_1 = 1,\;r_2 = 2$。

设 $a_n = A \cdot 1^n + B \cdot 2^n = A + B \cdot 2^n$。

由初始条件：
\[
\begin{cases}
a_1 = A + 2B = 1 \\
a_2 = A + 4B = 3
\end{cases}
\]
解得 $B = 1,\;A = -1$。

故 $a_n = 2^n - 1$。

\textbf{验证：}$a_1 = 1,\;a_2 = 3,\;a_3 = 3\times 3 - 2\times 1 = 7 = 2^3 - 1$。\;\checkmark

% ------ 例题2 ------
\begin{examplebox}[【例题2】重根情形的特征方程法]
已知 $a_1 = 1,\;a_2 = 3$，且 $a_{n+2} = 4a_{n+1} - 4a_n$，求 $a_n$。
\end{examplebox}

\textbf{解：}特征方程 $r^2 - 4r + 4 = 0$，即 $(r-2)^2 = 0$，重根 $r = 2$。

设 $a_n = (A + Bn) \cdot 2^n$。

由初始条件：
\[
\begin{cases}
a_1 = (A+B) \cdot 2 = 1 \\
a_2 = (A+2B) \cdot 4 = 3
\end{cases}
\implies
\begin{cases}
2A + 2B = 1 \\
4A + 8B = 3
\end{cases}
\]
解得 $A = \dfrac{1}{4},\;B = \dfrac{1}{4}$。

故 $a_n = \dfrac{n+1}{4} \cdot 2^n = (n+1)\cdot 2^{n-2}$。

% ------ 例题3 ------
\begin{examplebox}[【例题3】不动点法解分式线性递推]
已知 $a_1 = 2$，$a_{n+1} = \dfrac{2a_n - 1}{a_n}$，求 $a_n$。
\end{examplebox}

\textbf{解：}令 $x = \dfrac{2x-1}{x}$，则 $x^2 = 2x - 1$，即 $(x-1)^2 = 0$，不动点 $x_0 = 1$。

令 $b_n = a_n - 1$，则
\[
b_{n+1} + 1 = \frac{2(b_n+1)-1}{b_n+1} = \frac{2b_n + 1}{b_n + 1}
\]
\[
b_{n+1} = \frac{2b_n + 1}{b_n + 1} - 1 = \frac{b_n}{b_n + 1}
\]
取倒数：
\[
\frac{1}{b_{n+1}} = \frac{b_n + 1}{b_n} = 1 + \frac{1}{b_n}
\]
令 $c_n = \dfrac{1}{b_n}$，则 $c_{n+1} = c_n + 1$，为等差数列。

$a_1 = 2$，故 $b_1 = 1$，$c_1 = 1$，所以 $c_n = n$。

回代：$b_n = \dfrac{1}{n}$，$a_n = 1 + \dfrac{1}{n} = \dfrac{n+1}{n}$。

% ------ 例题4 ------
\begin{examplebox}[【例题4】双不动点分式递推]
已知 $a_1 = 3$，$a_{n+1} = \dfrac{a_n + 3}{a_n + 1}$，求 $a_n$。
\end{examplebox}

\textbf{解：}不动点方程 $x = \dfrac{x+3}{x+1}$，即 $x^2 + x = x + 3$，$x^2 = 3$。

故 $\alpha = \sqrt{3},\;\beta = -\sqrt{3}$。

构造：
\[
\frac{a_{n+1} - \sqrt{3}}{a_{n+1} + \sqrt{3}} = \frac{\frac{a_n+3}{a_n+1} - \sqrt{3}}{\frac{a_n+3}{a_n+1} + \sqrt{3}} = \frac{(a_n+3) - \sqrt{3}(a_n+1)}{(a_n+3) + \sqrt{3}(a_n+1)}
\]
\[
= \frac{a_n(1-\sqrt{3}) + (3-\sqrt{3})}{a_n(1+\sqrt{3}) + (3+\sqrt{3})} = \frac{(1-\sqrt{3})(a_n - \sqrt{3})}{(1+\sqrt{3})(a_n + \sqrt{3})}
\]
令 $t_n = \dfrac{a_n - \sqrt{3}}{a_n + \sqrt{3}}$，则 $t_{n+1} = \dfrac{1-\sqrt{3}}{1+\sqrt{3}} \cdot t_n$。

记 $\lambda = \dfrac{1-\sqrt{3}}{1+\sqrt{3}} = \dfrac{(1-\sqrt{3})^2}{1-3} = \dfrac{4-2\sqrt{3}}{-2} = \sqrt{3}-1$。

故 $t_n = t_1 \cdot (\sqrt{3}-1)^{n-1}$。

$a_1 = 3$，$t_1 = \dfrac{3-\sqrt{3}}{3+\sqrt{3}} = \dfrac{(3-\sqrt{3})^2}{9-3} = \dfrac{12-6\sqrt{3}}{6} = 2-\sqrt{3}$。

故 $\dfrac{a_n - \sqrt{3}}{a_n + \sqrt{3}} = (2-\sqrt{3})(\sqrt{3}-1)^{n-1}$。

解出：$a_n = \sqrt{3} \cdot \dfrac{1 + (2-\sqrt{3})(\sqrt{3}-1)^{n-1}}{1 - (2-\sqrt{3})(\sqrt{3}-1)^{n-1}}$。

% ------ 例题5 ------
\begin{examplebox}[【例题5】递推数列的周期性]
设 $a_1 = 2$，$a_{n+1} = \dfrac{1+a_n}{1-a_n}$，证明数列 $\{a_n\}$ 为周期数列，并求 $a_{2023}$。
\end{examplebox}

\textbf{解：}计算前几项：
\[
a_1 = 2,\quad a_2 = \frac{1+2}{1-2} = -3,\quad a_3 = \frac{1+(-3)}{1-(-3)} = \frac{-2}{4} = -\frac{1}{2}
\]
\[
a_4 = \frac{1+(-\frac{1}{2})}{1-(-\frac{1}{2})} = \frac{\frac{1}{2}}{\frac{3}{2}} = \frac{1}{3},\quad a_5 = \frac{1+\frac{1}{3}}{1-\frac{1}{3}} = \frac{\frac{4}{3}}{\frac{2}{3}} = 2
\]

$a_5 = a_1 = 2$，故周期 $T = 4$。

$a_{2023}$：$2023 = 4 \times 505 + 3$，故 $a_{2023} = a_3 = -\dfrac{1}{2}$。

\begin{tipbox}[方法总结]
对于分式线性递推 $a_{n+1} = \dfrac{a\,a_n+b}{c\,a_n+d}$，计算前若干项即可发现周期性。若不动点为虚数，则数列往往具有周期性。本题不动点方程 $x = \dfrac{1+x}{1-x}$ 即 $x - x^2 = 1 + x$，得 $x^2 = -1$，无实不动点，恰好对应周期4。
\end{tipbox}

% ------ 例题6 ------
\begin{examplebox}[【例题6】数学归纳法证明数列不等式]
已知 $a_1 = 1$，$a_{n+1} = a_n + \dfrac{1}{a_n}$，证明：$\sqrt{2n-1} \leq a_n \leq \sqrt{3n-2}$。
\end{examplebox}

\textbf{证明：}

\textbf{（下界）}先证 $a_n \geq \sqrt{2n-1}$。

由递推式 $a_{n+1}^2 = a_n^2 + 2 + \dfrac{1}{a_n^2} > a_n^2 + 2$。

累加：$a_n^2 > a_1^2 + 2(n-1) = 1 + 2n - 2 = 2n - 1$。

故 $a_n > \sqrt{2n-1}$。\;\checkmark

\textbf{（上界）}再证 $a_n \leq \sqrt{3n-2}$。

$a_{n+1}^2 = a_n^2 + 2 + \dfrac{1}{a_n^2}$。

由下界 $a_n^2 \geq 2n-1$，故 $\dfrac{1}{a_n^2} \leq \dfrac{1}{2n-1}$。

$a_{n+1}^2 \leq a_n^2 + 2 + \dfrac{1}{2n-1}$。

累加：$a_n^2 \leq 1 + 2(n-1) + \sum_{k=1}^{n-1}\dfrac{1}{2k-1}$。

需证 $\displaystyle\sum_{k=1}^{n-1}\dfrac{1}{2k-1} \leq n - 1$。

这等价于 $\displaystyle\sum_{k=1}^{n-1}\dfrac{1}{2k-1} \leq \sum_{k=1}^{n-1} 1$，即 $\dfrac{1}{2k-1} \leq 1$，对 $k\geq 1$ 显然成立。

故 $a_n^2 \leq 1 + 2(n-1) + (n-1) = 3n - 2$，即 $a_n \leq \sqrt{3n-2}$。\;\checkmark

% ------ 例题7 ------
\begin{examplebox}[【例题7】裂项求和与放缩]
证明：$\displaystyle\sum_{k=1}^{n} \frac{1}{k^2} < 2$。
\end{examplebox}

\textbf{证明：}$\dfrac{1}{k^2} < \dfrac{1}{k(k-1)} = \dfrac{1}{k-1} - \dfrac{1}{k}$（$k \geq 2$）。

\[
\sum_{k=1}^{n}\frac{1}{k^2} = 1 + \sum_{k=2}^{n}\frac{1}{k^2} < 1 + \sum_{k=2}^{n}\left(\frac{1}{k-1} - \frac{1}{k}\right) = 1 + 1 - \frac{1}{n} = 2 - \frac{1}{n} < 2
\]

\textbf{证毕。}

\begin{tipbox}[更强的上界]
利用 $\dfrac{1}{k^2} < \dfrac{1}{k^2 - \frac{1}{4}} = \dfrac{1}{(k-\frac{1}{2})(k+\frac{1}{2})} = \dfrac{1}{2k-1} - \dfrac{1}{2k+1}$，可以得到更强的结果 $\displaystyle\sum_{k=1}^{n}\dfrac{1}{k^2} < \dfrac{\pi^2}{6} \approx 1.6449$（但证明更复杂）。
\end{tipbox}

% ------ 例题8 ------
\begin{examplebox}[【例题8】裂项放缩求和的精确估计]
证明：$\displaystyle\frac{n}{2n+1} < \sum_{k=1}^{n}\frac{1}{(2k)^2 - 1} < \frac{n}{2n-1}$。
\end{examplebox}

\textbf{证明：}$\dfrac{1}{(2k)^2 - 1} = \dfrac{1}{4k^2 - 1} = \dfrac{1}{(2k-1)(2k+1)} = \dfrac{1}{2}\left(\dfrac{1}{2k-1} - \dfrac{1}{2k+1}\right)$。

\[
\sum_{k=1}^{n}\frac{1}{4k^2-1} = \frac{1}{2}\left(1 - \frac{1}{2n+1}\right) = \frac{n}{2n+1}
\]

等等——此处恰好等于 $\dfrac{n}{2n+1}$，等号成立，不等式左端应取等号。

重新检查：题目要求严格不等号。若确实要求证明严格不等式，则需修改估计方式。这里我们验证 $\dfrac{n}{2n+1} \leq \dfrac{n}{2n+1}$ 恰好取等。

实际上，令 $S_n = \displaystyle\sum_{k=1}^{n}\dfrac{1}{4k^2-1} = \dfrac{n}{2n+1}$。

验证右端：$\dfrac{n}{2n+1} < \dfrac{n}{2n-1}$，即 $2n-1 < 2n+1$，显然成立。

验证左端：$S_n = \dfrac{n}{2n+1}$，左端即为 $\dfrac{n}{2n+1}$，取等号。

故准确结论为：$\dfrac{n}{2n+1} \leq \displaystyle\sum_{k=1}^{n}\dfrac{1}{4k^2-1} < \dfrac{n}{2n-1}$，左端取等。

% ------ 例题9 ------
\begin{examplebox}[【例题9】构造辅助数列证明不等式]
已知 $a_1 = 1$，$a_{n+1} = \dfrac{a_n^2 + 2}{2a_n}$，证明：$a_n \geq \sqrt{2}$（$n\geq 2$），并求 $\{a_n\}$ 的极限。
\end{examplebox}

\textbf{解：}

\textbf{（1）证明下界：}$a_{n+1} = \dfrac{a_n^2 + 2}{2a_n} = \dfrac{a_n}{2} + \dfrac{1}{a_n}$。

由均值不等式（$a_n > 0$ 时）：
\[
a_{n+1} = \frac{a_n}{2} + \frac{1}{a_n} \geq 2\sqrt{\frac{a_n}{2}\cdot\frac{1}{a_n}} = 2\sqrt{\frac{1}{2}} = \sqrt{2}
\]

等号条件 $\dfrac{a_n}{2} = \dfrac{1}{a_n}$，即 $a_n = \sqrt{2}$。

$a_1 = 1$，$a_2 = \dfrac{1+2}{2} = \dfrac{3}{2} > \sqrt{2}$。

由归纳法，$n\geq 2$ 时 $a_n \geq \sqrt{2}$。\;\checkmark

\textbf{（2）求极限：}设 $\lim a_n = L$，则 $L = \dfrac{L^2+2}{2L}$，得 $2L^2 = L^2 + 2$，$L^2 = 2$，$L = \sqrt{2}$。

进一步验证单调性（$n\geq 2$）：
\[
a_{n+1} - a_n = \frac{a_n}{2} + \frac{1}{a_n} - a_n = \frac{1}{a_n} - \frac{a_n}{2} = \frac{2 - a_n^2}{2a_n}
\]
当 $a_n > \sqrt{2}$ 时 $a_{n+1} < a_n$（递减），当 $a_n = \sqrt{2}$ 时取等。故 $a_n$ 单调递减有下界 $\sqrt{2}$，极限为 $\sqrt{2}$。

% ------ 例题10 ------
\begin{examplebox}[【例题10】递推与不等式综合]
已知 $a_1 = \dfrac{1}{2}$，$a_{n+1} = a_n^2 + a_n$，证明：$\dfrac{1}{a_1 + 1} + \dfrac{1}{a_2 + 1} + \cdots + \dfrac{1}{a_n + 1} < 2$。
\end{examplebox}

\textbf{证明：}由 $a_{n+1} = a_n(a_n + 1)$，得 $\dfrac{1}{a_{n+1}} = \dfrac{1}{a_n(a_n+1)} = \dfrac{1}{a_n} - \dfrac{1}{a_n+1}$。

故 $\dfrac{1}{a_n + 1} = \dfrac{1}{a_n} - \dfrac{1}{a_{n+1}}$。

\[
\sum_{k=1}^{n}\frac{1}{a_k+1} = \sum_{k=1}^{n}\left(\frac{1}{a_k} - \frac{1}{a_{k+1}}\right) = \frac{1}{a_1} - \frac{1}{a_{n+1}} = 2 - \frac{1}{a_{n+1}}
\]

由 $a_1 = \dfrac{1}{2}$，$a_2 = \dfrac{1}{4} + \dfrac{1}{2} = \dfrac{3}{4}$，$a_n > 0$（由归纳法易证），故 $\dfrac{1}{a_{n+1}} > 0$。

\[
\sum_{k=1}^{n}\frac{1}{a_k+1} = 2 - \frac{1}{a_{n+1}} < 2
\]

\textbf{证毕。}

\begin{tipbox}[方法总结]
本题的关键是发现 $\dfrac{1}{a_n+1} = \dfrac{1}{a_n} - \dfrac{1}{a_{n+1}}$，即数列 $\left\{\dfrac{1}{a_n}\right\}$ 的差恰好等于目标求和项。这类"望远镜求和"在递推数列中非常常见。
\end{tipbox}

% ------ 例题11 ------
\begin{examplebox}[【例题11】特征方程法解决实际递推问题]
某实验室培养皿中的细菌数量 $a_n$（万个）满足 $a_1 = 1,\;a_2 = 5$，且 $a_{n+2} = 4a_{n+1} - 3a_n$。求 $a_n$ 并判断 $a_{10}$ 是否超过1000万个。
\end{examplebox}

\textbf{解：}特征方程 $r^2 - 4r + 3 = 0$，$(r-1)(r-3) = 0$，$r_1 = 1,\;r_2 = 3$。

设 $a_n = A + B \cdot 3^n$。

\[
\begin{cases}
A + 3B = 1 \\
A + 9B = 5
\end{cases}
\implies B = \frac{2}{3},\;A = -1
\]

$a_n = \dfrac{2}{3}\cdot 3^n - 1 = 2\cdot 3^{n-1} - 1$。

$a_{10} = 2\cdot 3^9 - 1 = 2\times 19683 - 1 = 39365$（万个）$= 3936.5$ 万个 $> 1000$ 万个。\;\checkmark

% ------ 例题12 ------
\begin{examplebox}[【例题12】数列不等式的归纳证明]
已知 $a_1 = 1$，$a_{n+1} = 1 + \dfrac{n}{a_n}$，证明：$\sqrt{n} \leq a_n \leq \sqrt{n} + 1$。
\end{examplebox}

\textbf{证明：}

\textbf{（下界 $a_n \geq \sqrt{n}$，用数学归纳法）}

$n=1$：$a_1 = 1 = \sqrt{1}$。\;\checkmark

设 $a_k \geq \sqrt{k}$，则
\[
a_{k+1} = 1 + \frac{k}{a_k} \leq 1 + \frac{k}{\sqrt{k}} = 1 + \sqrt{k}
\]
（注意这里得到的是上界估计，下界需要另一方向。）

$a_{k+1}^2 = \left(1+\dfrac{k}{a_k}\right)^2 = 1 + \dfrac{2k}{a_k} + \dfrac{k^2}{a_k^2}$。

由 $a_k \leq \sqrt{k}+1$（归纳假设上界）：
\[
a_{k+1}^2 \geq 1 + \frac{2k}{\sqrt{k}+1} + \frac{k^2}{(\sqrt{k}+1)^2} = \frac{(\sqrt{k}+1)^2 + 2k(\sqrt{k}+1) + k^2}{(\sqrt{k}+1)^2}
\]
分子 $= (\sqrt{k}+1+k)^2$，故 $a_{k+1} \geq \dfrac{\sqrt{k}+1+k}{\sqrt{k}+1} = 1 + \dfrac{k}{\sqrt{k}+1}$。

需证 $1 + \dfrac{k}{\sqrt{k}+1} \geq \sqrt{k+1}$。

$1 + \dfrac{k}{\sqrt{k}+1} = \dfrac{\sqrt{k}+1+k}{\sqrt{k}+1} = \dfrac{(\sqrt{k}+1)^2 - 1}{\sqrt{k}+1} = \sqrt{k}+1 - \dfrac{1}{\sqrt{k}+1}$。

需证 $\sqrt{k}+1 - \dfrac{1}{\sqrt{k}+1} \geq \sqrt{k+1}$，即 $1 - \dfrac{1}{\sqrt{k}+1} \geq \sqrt{k+1} - \sqrt{k} = \dfrac{1}{\sqrt{k+1}+\sqrt{k}}$。

即 $\dfrac{\sqrt{k}}{\sqrt{k}+1} \geq \dfrac{1}{\sqrt{k+1}+\sqrt{k}}$，即 $\sqrt{k}(\sqrt{k+1}+\sqrt{k}) \geq \sqrt{k}+1$。

左边 $= \sqrt{k}\sqrt{k+1} + k \geq k \geq \sqrt{k}+1$（$k\geq 2$ 时成立）。

$k=1$：$a_2 = 1+\frac{1}{1} = 2 \geq \sqrt{2}$。\;\checkmark

\textbf{（上界 $a_n \leq \sqrt{n}+1$，数学归纳法）}

$n=1$：$a_1 = 1 \leq 1+1 = 2$。\;\checkmark

设 $a_k \leq \sqrt{k}+1$，则
\[
a_{k+1} = 1 + \frac{k}{a_k} \leq 1 + \frac{k}{\sqrt{k}} = 1 + \sqrt{k} < 1 + \sqrt{k+1}
\]

\textbf{证毕。}

% ============ 第三部分：应用题 ============
\section*{\textcolor{apppurple}{三、应用题}}

\begin{appbox}[【应用1】贷款还款问题]
小张向银行贷款 $100$ 万元，年利率 $5\%$（复利），每年末还款 $x$ 万元。设第 $n$ 年初欠款为 $a_n$。

(1) 写出 $a_n$ 的递推关系；

(2) 若要求 $10$ 年还清，每年至少还款多少万元？

（参考：$1.05^{10} \approx 1.6289$）
\end{appbox}

\textbf{解：}

(1) 第 $n$ 年初欠款 $a_n$，年末变为 $1.05a_n$，还款 $x$ 后：
\[
a_{n+1} = 1.05a_n - x,\quad a_1 = 100
\]

(2) 这是一阶线性递推，不动点 $r = \dfrac{x}{1.05-1} = \dfrac{x}{0.05} = 20x$。

令 $b_n = a_n - 20x$，则 $b_{n+1} = 1.05b_n$，$b_n = (a_1 - 20x)\cdot 1.05^{n-1}$。

$a_n = 20x + (100 - 20x)\cdot 1.05^{n-1}$。

$10$ 年还清要求 $a_{11} \leq 0$：
\[
20x + (100 - 20x)\cdot 1.05^{10} \leq 0
\]
\[
20x \geq (20x - 100)\cdot 1.05^{10}
\]
\[
20x(1.05^{10} - 1) \leq 100 \cdot 1.05^{10}
\]
\[
x \leq \frac{100 \times 1.05^{10}}{20(1.05^{10}-1)} = \frac{100\times 1.6289}{20\times 0.6289} \approx \frac{162.89}{12.578} \approx 12.95
\]

注意这是最大还款额约束。由 $a_{11}\leq 0$ 得 $x \geq 12.95$（万元），即每年至少还款约 $12.95$ 万元。

\begin{appbox}[【应用2】传染病传播模型]
某传染病传播初期，每日新增感染者 $a_n$ 满足 $a_1 = 10$，$a_{n+1} = 1.2a_n - 5$。问：该模型下感染者最终会趋于稳定还是爆发？若稳定，每日新增多少人？
\end{appbox}

\textbf{解：}不动点方程 $r = 1.2r - 5$，解得 $r = 25$。

令 $b_n = a_n - 25$，则 $b_{n+1} = 1.2b_n$，$b_n = (a_1-25)\cdot 1.2^{n-1} = -15\cdot 1.2^{n-1}$。

$a_n = 25 - 15\cdot 1.2^{n-1}$。

由于 $1.2 > 1$，$15\cdot 1.2^{n-1}$ 递增，$a_n$ 会递减变为负数。

这意味着该模型下感染者会在某天后变为负数（不合理），说明模型只在初期有效。实际上当 $a_n$ 降至 $0$ 附近时，线性模型已不适用。

$a_n = 0$ 时 $25 = 15\cdot 1.2^{n-1}$，$1.2^{n-1} = \dfrac{5}{3}$，$(n-1)\ln 1.2 = \ln\dfrac{5}{3}$，$n \approx 1 + \dfrac{0.5108}{0.1823} \approx 3.8$，约第 $4$ 天感染者降至 $0$。

\begin{appbox}[【应用3】投资回报与数列不等式]
某基金投资方案：初始投入 $a_0 = 10$ 万元，第 $n$ 年末资产 $a_n = 1.08a_{n-1} + 2$（每年追加投入 $2$ 万元）。证明 $15$ 年后资产超过 $80$ 万元。
\end{appbox}

\textbf{解：}$a_n = 1.08a_{n-1} + 2$，不动点 $r = \dfrac{2}{1.08-1} = \dfrac{2}{0.08} = 25$。

令 $b_n = a_n - 25$，则 $b_n = 1.08b_{n-1}$，$b_n = (a_0 - 25)\cdot 1.08^n = -15\cdot 1.08^n$。

$a_n = 25 - 15\cdot 1.08^n$。

$a_{15} = 25 - 15\times 1.08^{15}$。

$1.08^{15} = e^{15\ln 1.08} \approx e^{15\times 0.07696} = e^{1.1544} \approx 3.1722$。

$a_{15} \approx 25 - 15\times 3.1722 = 25 - 47.58 = -22.58$？

这是负值，说明不动点 $r=25$ 是不稳定的（$1.08>1$ 使得 $b_n$ 发散）。实际上我们需要直接展开：

$a_n = 1.08^n \cdot a_0 + 2\cdot\dfrac{1.08^n - 1}{1.08 - 1} = 10\cdot 1.08^n + 25(1.08^n - 1) = 35\cdot 1.08^n - 25$。

$a_{15} = 35\times 1.08^{15} - 25 \approx 35\times 3.1722 - 25 = 111.03 - 25 = 86.03$（万元）$> 80$ 万元。\;\checkmark

% ============ 第四部分：课堂练习 ============
\section*{\textcolor{tipgreen}{四、课堂练习}}

\begin{enumerate}
\item 已知 $a_1 = 1,\;a_2 = 2$，$a_{n+2} = 5a_{n+1} - 6a_n$，求 $a_n$。

\item 已知 $a_1 = 2$，$a_{n+1} = \dfrac{3a_n + 2}{a_n + 2}$，用不动点法求 $a_n$。

\item 已知 $a_1 = 1$，$a_{n+1} = 2 - \dfrac{1}{a_n}$，证明 $\{a_n\}$ 为周期数列并求 $a_{100}$。

\item 已知 $a_n = \dfrac{1}{\sqrt{n} + \sqrt{n+1}}$，求 $S_n = \displaystyle\sum_{k=1}^{n}a_k$。

\item 证明：$\displaystyle\sum_{k=2}^{n}\frac{1}{k^2} < 1$（$n\geq 2$）。

\item 已知 $a_1 = 1$，$a_{n+1} = a_n + \dfrac{1}{a_n^2}$，证明 $a_n^3 \geq 3n - 2$。

\item 已知 $a_1 = \sqrt{2}$，$a_{n+1} = \dfrac{a_n^2 + 2}{2a_n}$，求 $\displaystyle\lim_{n\to\infty}a_n$ 并证明 $a_n \geq \sqrt{2}$。

\item 已知 $a_1 = 1$，$a_{n+1} = \dfrac{a_n}{1+a_n}$，求 $a_n$，并证明 $\displaystyle\sum_{k=1}^{n}a_k < 1$。
\end{enumerate}

\vspace{12pt}

\begin{knowledgebox}[参考答案]
\textbf{1.} 特征方程 $r^2 - 5r + 6 = 0$，$r_1 = 2,\;r_2 = 3$。

$a_n = A\cdot 2^n + B\cdot 3^n$，由 $a_1=1,\;a_2=2$：$2A+3B=1,\;4A+9B=2$，解得 $A = \dfrac{1}{2},\;B = 0$。

$a_n = 2^{n-1}$。

\vspace{6pt}
\textbf{2.} 不动点 $x = \dfrac{3x+2}{x+2}$，$x^2+2x=3x+2$，$x^2-x-2=0$，$x_1=2,\;x_2=-1$。

$\dfrac{a_n-2}{a_n+1}$：$\dfrac{a_{n+1}-2}{a_{n+1}+1} = \dfrac{\frac{3a_n+2}{a_n+2}-2}{\frac{3a_n+2}{a_n+2}+1} = \dfrac{a_n-2}{4a_n+4}$。

令 $t_n = \dfrac{a_n-2}{a_n+1}$，经计算 $t_{n+1} = \dfrac{1}{4}t_n$，$t_1 = 0$，故 $a_n = 2$（常数列）。

（验证：$a_2 = \dfrac{6+2}{2+2} = 2$。\;\checkmark）

\vspace{6pt}
\textbf{3.} 计算：$a_1=1,\;a_2=1,\;a_3=1,\ldots$ 所有项均为 $1$，$a_{100}=1$。

（验证：$a_2 = 2-1 = 1$，$a_3 = 2-1 = 1$。常数列，周期 $T=1$。）

\vspace{6pt}
\textbf{4.} $a_k = \dfrac{1}{\sqrt{k}+\sqrt{k+1}} = \sqrt{k+1}-\sqrt{k}$。

$S_n = \displaystyle\sum_{k=1}^{n}(\sqrt{k+1}-\sqrt{k}) = \sqrt{n+1}-1$。

\vspace{6pt}
\textbf{5.} $\dfrac{1}{k^2} < \dfrac{1}{k(k-1)} = \dfrac{1}{k-1}-\dfrac{1}{k}$。

$\displaystyle\sum_{k=2}^{n}\dfrac{1}{k^2} < \sum_{k=2}^{n}\left(\dfrac{1}{k-1}-\dfrac{1}{k}\right) = 1-\dfrac{1}{n} < 1$。\;\checkmark

\vspace{6pt}
\textbf{6.} $a_{n+1}^3 = \left(a_n+\dfrac{1}{a_n^2}\right)^3 = a_n^3 + 3 + \dfrac{3}{a_n^3} + \dfrac{1}{a_n^6} > a_n^3 + 3$。

累加：$a_n^3 > 1 + 3(n-1) = 3n-2$。\;\checkmark

\vspace{6pt}
\textbf{7.} 不动点 $x = \dfrac{x^2+2}{2x}$，$x^2=2$，$x=\sqrt{2}$（取正）。

$a_n - \sqrt{2} = \dfrac{(a_{n-1}-\sqrt{2})^2}{2a_{n-1}} \geq 0$（$a_{n-1}>0$ 时），故 $a_n \geq \sqrt{2}$。

$a_n$ 单调递减且有下界 $\sqrt{2}$，$\lim a_n = \sqrt{2}$。

\vspace{6pt}
\textbf{8.} $\dfrac{1}{a_{n+1}} = \dfrac{1+a_n}{a_n} = \dfrac{1}{a_n}+1$，故 $\dfrac{1}{a_n} = n$，$a_n = \dfrac{1}{n}$。

$S_n = \displaystyle\sum_{k=1}^{n}\dfrac{1}{k} < 1 + \displaystyle\sum_{k=2}^{n}\dfrac{1}{k(k-1)} = 1+1-\dfrac{1}{n} = 2-\dfrac{1}{n} < 2$。

但题目要求 $< 1$。实际上 $a_1 = 1$，$S_1 = 1$，$S_n = 1 + \displaystyle\sum_{k=2}^{n}\dfrac{1}{k}$，对 $n\geq 2$ 有 $S_n > 1$。

故 $\displaystyle\sum_{k=1}^{n}a_k < 1$ 不成立。正确结论应为 $\displaystyle\sum_{k=1}^{n}a_k < 2$，或 $\displaystyle\sum_{k=2}^{n}a_k < 1$。
\end{knowledgebox}

\vfill
\begin{center}
\small\textcolor{gray}{本讲义仅供教学参考，例题可根据学生水平调整难度。}
\end{center}

\end{document}
